\Askhsh{31530_SOLUTION}{1}
\begin{ylh}{1}
    \item [Ύλη από εικόνα]
\end{ylh}
\textbf{ΛΥΣΗ}
\begin{alist}
    \item[α)]
    \begin{rlist}
        \item Η συνάρτηση $f$ είναι παραγωγίσιμη με $f'(x)=3x^2+5$, $x\in\mathbb{R}$ και $f'(x)>0$, οπότε η $f$ είναι γνησίως αύξουσα στο $\mathbb{R}$.
        Επιπλέον, $f(0)=-2<0$ και $f(1)=1+5-2=4>0$, οπότε από το Θ. Bolzano υπάρχει $x_o \in (0, 1)$, μοναδικό λόγω της μονοτονίας της $f$, ώστε $f(x_o)=0$. Άρα η γραφική παράσταση της $f$ τέμνει τον άξονα $x'x$ σε ένα μόνο σημείο με τετμημένη $x_o$ που περιέχεται στο διάστημα $(0, 1)$.

        \item Ισχύει: $f\left(\frac{1}{2}\right) = \frac{1}{8}+\frac{5}{2}-\frac{1}{2} = \frac{1+20-4}{8} = \frac{17}{8} > 0$ και $f(0)=-2<0$, οπότε $f(0)\cdot f\left(\frac{1}{2}\right) < 0$. Επομένως το διάστημα στο οποίο περιέχεται το $x_o$ περιορίζεται στο $\left(0, \frac{1}{2}\right)$ οπότε ο αριθμός $x_o$ είναι πιο κοντά στο $0$ παρά στο $1$.
    \end{rlist}

    \item[β)] Ο αριθμός $\theta$ είναι θετικός, οπότε έχουμε:
    \[ x_o-\theta < x_o < x_o+\theta \xrightarrow{f\uparrow} f(x_o-\theta) < f(x_o) < f(x_o+\theta) \Rightarrow f(x_o-\theta) < 0 < f(x_o+\theta) \]
    απ' όπου προκύπτει ότι: $\frac{f(x_o+\theta)}{f(x_o-\theta)} < 0$, (1).
    Είναι:
    \[ \lim_{x\to+\infty} \frac{f(x_o+\theta)x^3+2x-5}{f(x_o-\theta)x-5} = \lim_{x\to+\infty} \frac{f(x_o+\theta)x^3}{f(x_o-\theta)x} = \frac{f(x_o+\theta)}{f(x_o-\theta)} \cdot \lim_{x\to+\infty} x^2 = -\infty \]
    λόγω της (1).

    \item[γ)] Η εξίσωση της εφαπτομένης της $C_f$ στο σημείο $A$ είναι:
    \[ y-f(1)=f'(1)(x-1) \Rightarrow y-4=8(x-1) \Rightarrow y=8x-4 \]
    Η συνάρτηση $f$ είναι συνεχής στο $[1,2]$ και ισχύει $f''(x)=6x>0$, για κάθε $x\in(1,2)$. Άρα η $f$ είναι κυρτή, οπότε η $C_f$ είναι πάνω από την εφαπτομένη της, με εξαίρεση το σημείο επαφής τους $A(1, 4)$. Έτσι, το ζητούμενο εμβαδόν $E$ είναι:
    \[ E = \int_1^2 [f(x)-(8x-4)]dx = \int_1^2 (x^3-3x+2)dx = \left[\frac{x^4}{4}-3\frac{x^2}{2}+2x\right]_1^2 \]
    \[ = 4-3\cdot 2+4-\left(\frac{1}{4}-\frac{3}{2}+2\right)=\frac{5}{4} \]
\end{alist}